\section{The hypothesis moduli from primitive data}\label{S:moduli}

The constants of Theorem~3.2 are not free assumptions; under the standing bounds they are computed from the data. Let $\|A_\theta\|_\infty\le\bar A$, $\|B_\theta\|_\infty\le\bar B$, $\|E_\theta\|_\infty\le\bar E$, $\|g_\theta\|_\infty\le\bar g$, $\|c_i\|\le\bar c$ on $[0,T]$; $U\subset c_U+R_U\mathbb B$, $D_\theta\subset R_D\mathbb B$, $P_\theta\subset R_P\mathbb B$; $\TV(B_\theta)\le V_B$, $\TV(E_\theta)\le V_E$; and $M_X:=e^{\bar AT}$, so that $\|X_\theta(t,\sigma)\|\le M_X$ (Gr\"onwall on the transition matrix) and $\partial_\sigma X_\theta(t,\sigma)=-X_\theta(t,\sigma)A_\theta(\sigma)$.

\paragraph{Total variation of the kernels (H2), (H4).}
For $k_a(\sigma)=\mathbf 1_{\{\sigma\le t\}}B_\theta(\sigma)^\top X_\theta(t,\sigma)^\top c_i$, the product rule for functions of bounded variation gives on $[0,t]$
\begin{multline*}
\TV(k_a;[0,t])\le\sup_\sigma\|X_\theta(t,\sigma)^\top c_i\|\,\TV(B_\theta)+\sup_\sigma\|B_\theta(\sigma)\|\int_0^t\|X_\theta(t,\sigma)A_\theta(\sigma)^\top c_i\|\dd\sigma\\
\le\bar c\,(M_XV_B+\bar BM_X\bar AT),
\end{multline*}
and the drop to zero at $\sigma=t$ adds at most $\|B_\theta(t)^\top c_i\|\le\bar B\bar c$. Hence $V_U=\bar c(\bar B+M_XV_B+\bar BM_X\bar AT)$, and identically $V_D=\bar c(\bar E+M_XV_E+\bar EM_X\bar AT)$ for the disturbance kernels.

\paragraph{$L^1$ label-continuity modulus (H3).}
For $t\ge t'$, Gr\"onwall on $X_\theta(t,\sigma)-X_\theta(t',\sigma)=\int_{t'}^tA_\theta(r)X_\theta(r,\sigma)\dd r$ gives $\|X_\theta(t,\sigma)-X_\theta(t',\sigma)\|\le M_X\bar A|t-t'|$ uniformly in $\sigma$, so
\[
\|k_t-k_{t'}\|_{L^1}\le\bar B\bar c\,|t-t'|+\int_0^{t'}\bar B\bar cM_X\bar A|t-t'|\dd\sigma\le\bar B\bar c(1+M_X\bar AT)|t-t'| ,
\]
and $|\beta_t-\beta_{t'}|\le L^{\mathrm{tr}}|t-t'|$ by the trajectory bound below. The label map is $L^1$-Lipschitz in $t$ with constant $\bar B\bar c(1+M_X\bar AT)+L^{\mathrm{tr}}$, which is the form the proof of Theorem~3.2(iv) consumes; the sup-norm form is false, because the step kernels jump at $\sigma=t$ for every $t$.

\paragraph{Uniform time-Lipschitz constant (H1).}
For a trajectory $x$ of mode $\theta$, $\|x(t)\|\le M_X\bigl(R_P+T(\bar g+\bar B(R_U+\|c_U\|)+\bar ER_D)\bigr)$ by Gr\"onwall, and
\begin{gather*}
\bigl|\tfrac{\mathrm d}{\mathrm dt}c_i^\top x(t)\bigr|\le\bar c\bigl(\bar A\|x(t)\|+\bar g+\bar B(R_U+\|c_U\|)+\bar ER_D\bigr)\le L^{\mathrm{tr}},\\
L^{\mathrm{tr}}:=\bar c\bar AM_XR_P+\bar c(1+\bar AM_XT)\bigl(\bar g+\bar B(R_U+\|c_U\|)+\bar ER_D\bigr).
\end{gather*}
Each $\calF_{(\theta,i,t)}(\pi)$ is a supremum over $(x_0,d)$ of $L^{\mathrm{tr}}$-Lipschitz functions of $t$, and a pointwise supremum of $L$-Lipschitz functions is $L$-Lipschitz; hence (H1) holds with $L=L^{\mathrm{tr}}$.

\paragraph{The instance.}
For the double integrator with $n=4$, $\bar A=1$, $\bar B=1$, $R_U=1$, $R_P=\tfr{34}{25}$ the closed forms apply but are loose; the direct kernel computation gives the sharp constants used in Section~5.5 of the main text. Position rows have kernels $k(\sigma)=(t-\sigma)n_j\mathbf 1_{\{\sigma\le t\}}$, continuous at $\sigma=t$, with $\TV=t\le T$; velocity rows have kernels $\pm e_i\mathbf 1_{\{\sigma\le t\}}$ with $\TV=1$: $V_U=\max(T,1)=\tfr65$. Since $|v_j(t)|\le1+t\le\tfr{11}5$ along every admissible trajectory and $|\dot v|\le1$, $L=\tfr{11}5$. Scenarios are exact (singleton $P_\theta$, no disturbance): $\delta_\beta=0$. For nonlinear or interval-enclosed data a validated flow-error bound enters at the same two places---it inflates the stored labels $\beta^+_a$ (hence $\delta_\beta$) and the kernel errors $\ebar_a$---and the sandwich (11) is then charged with the inflated budgets; the relaxation direction must be toward a larger feasible controller set, as in Section~8.1 of the main text.

\paragraph{The inner-label sharpening (9) and the bound $Th/4$.}
On a mesh whose label times are block boundaries, a block $J=[s_0,s_1]$ with midpoint $m$ lying before the label time $t$ of a position row has $\bar k-k=(\sigma-m)n_j$ on $J$, and
\[
\int_Jh_U\bigl((\sigma-m)n_j\bigr)\dd\sigma=\int_J|\sigma-m|\,h_U(\pm n_j)\dd\sigma=\tfrac14|J|^2 ,
\]
because $h_U(\pm n_j)=1$. Blocks after $t$ contribute nothing (the kernel vanishes), and velocity rows have piecewise-constant kernels with zero block error. Hence the exact moment error of the label $(j,\text{critical},t)$ is $e_t=\tfrac14\sum_{J\subseteq[0,t]}|J|^2$, which is $th/4\le Th/4$ on the uniform mesh and equals $\tfrac14\sum_J|J|^2\le Th/4$ at the horizon on any aligned mesh with $\max_J|J|=h$. The uniform bound $R_UV_Uh=\tfr65h$ is a factor $4$ larger at the horizon ($Th/4=\tfr3{10}h$). Both are legitimate inputs to program (10); the main text uses the sharp one.

\paragraph{The loose-budget example of Section~5.5.}
Feeding program (10) at $h=\tfr1{10}$ the inflated budget $\ebar_a=e_a+\tfr12$ on every row returns $\rho=\tfr3{100}-\tfr12=-\tfr{47}{100}$ (the dual witness is unchanged; the objective drops by the added budget), and the sandwich (11) gives $J\le-\tfr{47}{100}+\max_a\ebar_a+Lh_t=-\tfr{47}{100}+\tfr{53}{100}+\tfr{11}{100}=\tfr{17}{100}$, still valid. The formula with the uniform special case $R_UV_Uh_u$ in place of $\max_a\ebar_a$ would give $-\tfr{47}{100}+\tfr{12}{100}+\tfr{11}{100}=-\tfr6{25}<\tfr3{50}=J$, which is false: the sandwich must be charged with the budget the program actually used.
